\chapter{Messaging and Event-Driven Architecture}\label{ch:pl:messaging-and-event-driven-architecture}

A consumer of the fills topic crashes after applying a batch of messages to the positions and before recording how far it has read. When it
restarts it reads the batch again and applies it again, and the firm's position in one stock is double until someone compares it with the
drop copy. Nothing was broken: the messaging system delivered every message at least once, which is exactly what it promised. This chapter
is about what messaging systems promise, what they cannot, and how a consumer turns \omterm{def:pl:messaging-and-event-driven-architecture:delivery}{at-least-once delivery} into exactly-once effects --- with a
\omterm{def:pl:messaging-and-event-driven-architecture:log}{partitioned log} built in process, a \omterm{def:pl:position-and-pnl-services:service}{position service} run under a thousand crashes in three modes, and the outbox that lets a service change
its database and announce it without losing either.

\section{Why systems talk through messages}

The platform map of chapter~\ref{ch:pl:the-platform-map} is a set of services: \omterm{def:pl:trade-capture-and-booking:capture}{trade capture}, positions, risk, P\&L, surveillance. They can
call each other directly, each one knowing who needs what; or each can publish what happened, and whoever needs it subscribes. The second
decouples them in three ways --- in knowledge (the publisher does not know its consumers), in time (a consumer can be down and catch up), and
in rate (a slow consumer does not slow the publisher) --- and makes the flow of events itself a record that can be replayed.

\begin{definition}[Publish--subscribe]\label{def:pl:messaging-and-event-driven-architecture:pubsub}
\emph{Publish--subscribe}\index{publish--subscribe} is a messaging pattern in which producers publish messages to named topics without
addressing any consumer, and each consumer receives the messages of the topics it subscribes to.
\end{definition}

The cost is that the call graph becomes invisible, and failures move from the call (which fails loudly) to the delivery (which can lose or
repeat messages quietly). Most of this chapter is about the second.

\section{Brokers and brokerless transports}

\begin{definition}[Message broker, brokerless messaging]\label{def:pl:messaging-and-event-driven-architecture:broker}
A \emph{message broker}\index{message broker} is a server that receives messages from producers, stores them, and delivers them to consumers,
so that producers and consumers only know the broker. In \emph{brokerless messaging}\index{brokerless messaging} producers send directly to
consumers --- by unicast, multicast or shared memory --- through a library in each process, with no server in the path.
\end{definition}

The two answer different needs. A broker adds a hop and a server to run, and gives in return storage, replay, and consumers that come and go.
Brokerless transports give the lowest latency and no central component: the market-data and order paths of Book~13 are brokerless, with
multicast (Book~13, chapter~16) and ring buffers in shared memory, and libraries such as Aeron --- \textquotedbl{}efficient reliable UDP
unicast, UDP multicast, and IPC message transport\textquotedbl{} --- provide the reliability on top. Between services that care more about
not losing an event than about microseconds --- \omterm{def:pl:trade-capture-and-booking:capture}{trade capture}, positions, risk, the post-trade chain --- the firm uses a log.

\section{The log as the backbone}

\begin{definition}[Partitioned log, consumer offset]\label{def:pl:messaging-and-event-driven-architecture:log}
A \emph{partitioned log}\index{partitioned log} stores each topic as several partitions, each an append-only sequence of messages addressed by
their position in it; a message's partition is chosen from its key, so that messages with the same key keep their order. A \emph{consumer
offset}\index{consumer offset} is the position up to which a consumer has read a partition, recorded by the consumer (committed) so that it
can resume there.
\end{definition}

\begin{definition}[Consumer group]\label{def:pl:messaging-and-event-driven-architecture:group}
A \emph{consumer group}\index{consumer group} is a set of consumer processes that share a topic's partitions, each partition read by exactly
one member at a time, with one committed offset per partition for the whole group; when members join or leave, the partitions are reassigned
(a rebalance).
\end{definition}

This is the design Kreps, Narkhede and Rao described for Kafka in 2011: a message has no identifier of its own, \textquotedbl{}each message is
addressed by its logical offset in the log\textquotedbl{}. The chapter's \texttt{firm.eventlog} builds it in process: topics, partitions by a
stable hash of the key, append-only segments on disk, retention that drops whole segments, and committed offsets per group. Keying fills by
account puts all of an account's fills in one partition, in order; eight partitions shared by three consumers are assigned three, three and
two.

\begin{omfigure}
\begin{tikzpicture}[font=\small,
  cell/.style={draw=omDef,fill=omDef!8,minimum width=6mm,minimum height=6mm,inner sep=0pt,font=\scriptsize},
  arr/.style={-{Stealth[length=2mm]},thick}]
\node[draw=omDef,rounded corners=2pt,align=center,text width=16mm,minimum height=24mm] (prod) at (-0.9,0) {producers\\(trade\\capture)};
\foreach \p/\n in {0/7,1/5,2/8,3/6} {
  \node[font=\scriptsize,anchor=east] at (2.4,1.35-0.9*\p) {partition \p};
  \foreach \i in {1,...,\n} \node[cell] at (2.3+0.62*\i,1.35-0.9*\p) {\the\numexpr\i-1\relax};
}
\draw[arr] (prod.east) -- (0.8,0);
\draw[omThm,very thick] (2.3+0.62*4.5,1.7) -- (2.3+0.62*4.5,1.0) node[above,pos=0,font=\scriptsize]{committed 4};
\draw[omThm,very thick] (2.3+0.62*3.5,-1.0) -- (2.3+0.62*3.5,-1.7) node[below,pos=1,font=\scriptsize]{committed 3};
\node[draw=omThm,rounded corners=2pt,align=center,text width=24mm] (c1) at (9.8,0.9) {consumer 1\\partitions 0, 1};
\node[draw=omThm,rounded corners=2pt,align=center,text width=24mm] (c2) at (9.8,-0.9) {consumer 2\\partitions 2, 3};
\node[font=\scriptsize,omThm] at (9.8,-1.9) {one consumer group};
\draw[arr] (6.95,1.35) -- (c1.west); \draw[arr] (5.7,0.45) -- (c1.west);
\draw[arr] (7.6,-0.45) -- (c2.west); \draw[arr] (6.35,-1.35) -- (c2.west);
\end{tikzpicture}

\omcaption{A topic of four partitions, each an append-only sequence of messages addressed by offset, and a \omterm{def:pl:messaging-and-event-driven-architecture:group}{consumer group} of two members that
share them. The group commits one offset per partition (red): the position from which a member resumes after a restart or a
rebalance.}\label{fig:pl:messaging-and-event-driven-architecture:log}
\end{omfigure}

\omcode{code/firm/eventlog/firm_eventlog.py}{53}{64}{Appending to a partitioned log: the key chooses the partition, and a producer's
(identifier, sequence) pair already stored is ignored, so that a resend does not duplicate.}\label{lst:pl:messaging-and-event-driven-architecture:append}

\section{Delivery guarantees and the exactly-once myth}

A consumer does two things with a batch: it processes the messages, and it commits the offset. A crash can fall between them, and the order
of the two decides what the crash costs.

\begin{definition}[At-most-once and at-least-once delivery]\label{def:pl:messaging-and-event-driven-architecture:delivery}
Under \emph{at-most-once delivery}\index{at-most-once delivery} a message may be lost but is never processed twice: the consumer commits the
offset before processing. Under \emph{at-least-once delivery}\index{at-least-once delivery} a message is never lost but may be processed more
than once: the consumer commits after processing, and a crash in between repeats the uncommitted messages.
\end{definition}

\begin{definition}[Exactly-once processing]\label{def:pl:messaging-and-event-driven-architecture:exactly}
\emph{Exactly-once processing}\index{exactly-once processing} is the guarantee that each message's effect on the consumer's state happens
once, whatever the crashes: obtained not from the transport, which can only deliver at least once, but by making the effect idempotent
(Book~13, chapter~24) or by storing the state and the offset in one atomic transaction.
\end{definition}

The 2011 paper was explicit: \textquotedbl{}Kafka only guarantees \omterm{def:pl:messaging-and-event-driven-architecture:delivery}{at-least-once delivery}\textquotedbl{}, and an application that cares about
duplicates \textquotedbl{}must add its own deduplication logic\textquotedbl{}. Later versions added an idempotent producer and transactions,
which give exactly-once between Kafka topics; for any other destination, the documentation says, exactly-once \textquotedbl{}generally
requires cooperation with such systems\textquotedbl{}. The \omterm{def:pl:position-and-pnl-services:service}{position service} is such a destination.

\omcode{code/firm/eventlog/firm_eventlog.py}{105}{126}{One batch under the three modes: commit before processing, commit after, or apply the
batch and the next offset in one database transaction.}\label{lst:pl:messaging-and-event-driven-architecture:modes}

The chapter's day is 10\,000 synthetic fills for ten accounts and twenty stocks, keyed by account into four partitions and consumed in batches
of 100. Each day ten crashes strike at points drawn uniformly over the consumer's steps --- every message processed, and the moment between
processing a batch and committing it --- and a hundred days are run in each mode. For every fill we count how many times it reached the
positions.

\begin{omfigure}
\begin{tikzpicture}
\begin{axis}[width=11.5cm,height=5.4cm,ybar,bar width=14pt,ylabel={fills per 1\,000 crashes},symbolic x coords={at-most-once,at-least-once,atomic},
  xtick=data,xticklabels={commit first\\(at most once),commit after\\(at least once),state and offset\\atomic},
  xticklabel style={align=center},
  grid=major,grid style={gray!20},tick label style={font=\small},label style={font=\small},table/col sep=comma,ymin=0,ymax=62000,
  nodes near coords,every node near coord/.append style={font=\scriptsize,/pgf/number format/1000 sep={\,}},scaled y ticks=false,
  yticklabel style={/pgf/number format/1000 sep={\,}},
  legend style={font=\small,at={(0.5,1.03)},anchor=south,legend columns=2,draw=none,/tikz/every even column/.append style={column sep=4mm}}]
\addplot[fill=omThm!60,draw=omThm] table[x=mode,y=lost_per_1000]{figdata/platforms/24-messaging-and-event-driven-architecture/modes.csv};
\addplot[fill=omDef!40,draw=omDef] table[x=mode,y=dup_per_1000]{figdata/platforms/24-messaging-and-event-driven-architecture/modes.csv};
\legend{lost,duplicated}
\end{axis}
\end{tikzpicture}

\omcaption{Fills lost and duplicated per 1\,000 consumer crashes, over a hundred days of ten crashes in each mode. Committing first loses about
half a batch per crash (50\,323 per 1\,000); committing after duplicates as much (48\,652); state and offset in one transaction, none. Data:
\texttt{fig\_eventlog.py}.}\label{fig:pl:messaging-and-event-driven-architecture:modes}
\end{omfigure}

The first two modes are mirror images (\cref{fig:pl:messaging-and-event-driven-architecture:modes}). A crash inside a batch of 100 falls on
average half-way, so committing first loses about 50 fills and committing after repeats about 50; a crash just after processing and before the
commit --- the hook's --- repeats all 100. In the third mode the positions and the next offset are written in one SQLite transaction: a crash
anywhere before the commit rolls both back, the consumer resumes from the offset that matches the state, and no fill is lost or counted
twice in 1\,000 crashes. The same effect comes from idempotent processing: store each fill's identifier with the state and skip identifiers
already applied.

\begin{omfigure}
\begin{tikzpicture}
\begin{axis}[width=11.5cm,height=5.4cm,xlabel={day, sorted by its error},ylabel={largest position error (shares)},
  grid=major,grid style={gray!20},tick label style={font=\small},label style={font=\small},table/col sep=comma,
  xmin=1,xmax=100,ymin=-200,ymax=7000,yticklabel style={/pgf/number format/1000 sep={\,}}]
\addplot[omThm,thick] table[x=rank,y=at_most_once]{figdata/platforms/24-messaging-and-event-driven-architecture/worst.csv};
\addplot[omDef,thick] table[x=rank,y=at_least_once]{figdata/platforms/24-messaging-and-event-driven-architecture/worst.csv};
\addplot[black,thick] table[x=rank,y=atomic]{figdata/platforms/24-messaging-and-event-driven-architecture/worst.csv};
\node[omDef,font=\scriptsize,anchor=east] at (axis cs:88,6300) {commit after: up to 6\,500};
\node[omThm,font=\scriptsize,anchor=north west] at (axis cs:55,3200) {commit first: up to 5\,100};
\node[font=\scriptsize,anchor=south west] at (axis cs:3,150) {atomic: 0 every day};
\end{axis}
\end{tikzpicture}

\omcaption{The largest end-of-day position error in any account and stock, for each of the hundred days, sorted. Ten crashes a day leave a
median worst error of 3\,500 shares when committing first and 3\,800 when committing after, and up to 5\,100 and 6\,500; the atomic mode ends
every day exact. Data: \texttt{fig\_eventlog.py}.}\label{fig:pl:messaging-and-event-driven-architecture:worst}
\end{omfigure}

\section{Patterns: outbox, event sourcing, dead letters}

The consumer's problem has a mirror on the producer's side. The trade-capture service must store a trade in its database and publish the
event that announces it. Two separate writes cannot be atomic: write then publish, and a crash in between leaves a trade nobody hears about;
publish then write, and consumers act on a trade that does not exist.

\begin{definition}[Transactional outbox]\label{def:pl:messaging-and-event-driven-architecture:outbox}
A \emph{transactional outbox}\index{transactional outbox} is a table in a service's own database into which the service writes each event in
the same transaction as the change it announces; a separate relay reads the table, publishes the events in order, and marks them sent, so
that an event is published if and only if its change was committed --- at least once, with a key that lets the log or the consumer drop the
repeats.
\end{definition}

\begin{omfigure}
\begin{tikzpicture}[font=\small,
  box/.style={draw=omDef,fill=omDef!8,rounded corners=2pt,align=center,minimum height=9mm,text width=24mm},
  arr/.style={-{Stealth[length=2mm]},thick}]
\node[box] (svc) {trade-capture\\service};
\node[box,right=12mm of svc,yshift=6mm] (t) {table \texttt{trades}};
\node[box,right=12mm of svc,yshift=-6mm] (o) {table \texttt{outbox}};
\node[draw=omThm,dashed,rounded corners=3pt,fit=(t)(o),inner sep=4pt,label={[font=\scriptsize,omThm]above:one transaction}] (tx) {};
\node[box,right=12mm of tx] (r) {relay};
\node[box,right=18mm of r] (l) {log, topic\\\texttt{trades}};
\draw[arr] (svc.east) -- (tx.west);
\draw[arr] (o.east) -- node[below,font=\scriptsize]{unsent} (r.west |- o.east);
\draw[arr] (r) -- node[above,font=\scriptsize]{1. publish} (l);
\draw[arr,omThm] (r.south) to[bend left=35] node[below,font=\scriptsize]{2. mark sent} (o.south east);
\end{tikzpicture}

\omcaption{The \omterm{def:pl:messaging-and-event-driven-architecture:outbox}{transactional outbox}: the trade and the event that announces it are committed in one database transaction; a relay publishes
unsent events to the log and then marks them sent. A crash between 1 and 2 repeats an event, never loses one, and the idempotence key lets
the log drop the repeat.}\label{fig:pl:messaging-and-event-driven-architecture:outbox}
\end{omfigure}

\omcode{code/firm/eventlog/firm_eventlog.py}{171}{192}{The outbox: the trade and its event in one transaction, and the relay that publishes
each event with its outbox number as idempotence key before marking it sent.}\label{lst:pl:messaging-and-event-driven-architecture:outbox}

\begin{table}[htbp]
\centering\small
\begin{tabular}{lrrr}
\toprule
design & events lost & phantom events & duplicated events\\
\midrule
write the trade, then publish & 22 & 0 & 0\\
publish, then write the trade & 0 & 22 & 0\\
outbox, relay without idempotence key & 0 & 0 & 22\\
outbox, relay with idempotence key & 0 & 0 & 0\\
\bottomrule
\end{tabular}
\caption{2\,000 trades published with a 1\% chance of a crash at the dangerous point of each design: between the two writes, or between
publishing and marking an outbox row sent. A phantom event announces a trade the database does not hold.}
\label{tab:pl:messaging-and-event-driven-architecture:outbox}
\end{table}

\Cref{tab:pl:messaging-and-event-driven-architecture:outbox} shows the three failure modes and the fix. The outbox moves the problem from
two systems to one, and what is left --- a relay that published and died before marking --- is a repeat, which the idempotent producer or
the consumer removes.

\begin{definition}[Event sourcing]\label{def:pl:messaging-and-event-driven-architecture:sourcing}
\emph{Event sourcing}\index{event sourcing} stores a system's state as the sequence of events that changed it, appended and never edited, and
derives every current state, and every past one, by replaying them.
\end{definition}

Chapter~\ref{ch:pl:trade-capture-and-booking}'s trade store is event-sourced, and the log is its natural home: a new downstream system
subscribes and replays the topic from its start to build its own state. The price is retention --- the log must keep every event, or a
snapshot with the offset it covers --- and discipline about changing the events' format.

\begin{definition}[Dead-letter queue]\label{def:pl:messaging-and-event-driven-architecture:dlq}
A \emph{dead-letter queue}\index{dead-letter queue} is a topic to which a consumer sends a message it has failed to process a stated number of
times, with the partition and offset it came from, so that one malformed message does not stop the partition behind it.
\end{definition}

In a \omterm{def:pl:messaging-and-event-driven-architecture:log}{partitioned log} a message that always fails is worse than in a queue: the consumer cannot skip it without committing past it, so every
message behind it in the partition waits. The dead-letter topic lets the consumer move on and leaves a record to be repaired and replayed ---
which, for fills, must be done before anyone trusts the positions.

\begin{dated}{2026-09}{Kafka's delivery semantics}\label{dat:pl:messaging-and-event-driven-architecture:kafka}
Kafka's documentation distinguishes at most once (messages may be lost but are never redelivered), at least once (never lost but may be
redelivered) and exactly once; a consumer that saves its position before processing gets the first, after processing the second. Since version
0.11.0.0 its producer offers an idempotent delivery option, so that resending does not create duplicate entries in the log, and transactions
that write to several partitions atomically; exactly-once delivery to other systems generally requires their cooperation.
\end{dated}

\section{Tutorial: a thousand crashes}\label{tut:pl:messaging-and-event-driven-architecture:tut}

\begin{tutorial}
\textbf{Goal.} Consume a day of fills under crashes in three modes and publish trades three ways. \textbf{End state:}
\cref{fig:pl:messaging-and-event-driven-architecture:modes} and \cref{tab:pl:messaging-and-event-driven-architecture:outbox}.
\begin{tutsteps}
\item \textbf{Log}: \texttt{firm\_eventlog.Log}, a topic of four partitions; \texttt{pl\_eventlog.publish(fills())}.
\item \textbf{One day}: \texttt{day(fs, mode, crash\_steps)} for each mode, with the same crash steps; count applications per fill.
\item \textbf{A hundred days}: \texttt{schedules(mode)}, lost and duplicated fills per 1\,000 crashes and the worst position error.
\item \textbf{Outbox}: \texttt{publishing(design)} for the four designs.
\item \textbf{Dead letters}: put a malformed fill in a partition and consume it with \texttt{process\_with\_dlq}.
\end{tutsteps}
\textbf{What to change next.} Replace the atomic store by idempotent processing (a table of applied fill identifiers) and check that it gives
the same result; rebalance the group in the middle of the day.
\end{tutorial}

\section{Build: the event log}\label{bld:pl:messaging-and-event-driven-architecture:eventlog}

\begin{build}
\bfield{Purpose} Carry the firm's events between services in order per key, durably, with consumers that can turn \omterm{def:pl:messaging-and-event-driven-architecture:delivery}{at-least-once delivery}
into exactly-once effects.
\bfield{Interface} \texttt{Log} (\texttt{create}, \texttt{append}, \texttt{read}, \texttt{end}, \texttt{retain}, \texttt{commit},
\texttt{committed}), \texttt{assign}, \texttt{run\_batch}, \texttt{AtomicStore}, \texttt{Outbox}, \texttt{process\_with\_dlq}.
\bfield{Rules} Partition by key; never edit a message; every consumer states its mode; state and offset atomic, or processing idempotent,
for any consumer whose state matters; every event published through an outbox; a message that keeps failing goes to dead letters with its
origin.
\bfield{Acceptance tests} \texttt{code/firm/eventlog/tests/}: one key in one partition and in order, segments and retention; a resend
ignored; the group assignment; the three modes under a crash; the outbox relay's resend deduplicated; a poison message to dead letters.
\bfield{Stretch} Consumer-group rebalancing with a generation number that fences out a stale member; compaction by key; a network server.
\end{build}

\begin{omsources}
\item J.~Kreps, N.~Narkhede and J.~Rao, \emph{Kafka: a Distributed Messaging System for Log Processing}, NetDB 2011.
\item Apache Kafka documentation, \emph{Design}: message delivery semantics.
\item One Quant Book~13, chapters~12, 16 and~24 (ring buffers and slow consumers, multicast, idempotent processing and sequencers).
\end{omsources}

\section{Exercises}

\begin{exercise}[$\star$]\label{exo:pl:messaging-and-event-driven-architecture:1}
Why are fills keyed by account rather than sent round-robin to the partitions?
\end{exercise}

\begin{exercise}[$\star$]\label{exo:pl:messaging-and-event-driven-architecture:2}
A consumer commits after processing and crashes after applying 30 messages of a batch of 100. What happens when it restarts?
\end{exercise}

\begin{exercise}[$\star$]\label{exo:pl:messaging-and-event-driven-architecture:3}
Eight partitions and three consumers in a group: how are they assigned, and what happens if a fourth joins?
\end{exercise}

\begin{exercise}[$\star\star$]\label{exo:pl:messaging-and-event-driven-architecture:4}
Why does committing first lose about 50 fills per crash, and committing after duplicate about as many?
\end{exercise}

\begin{exercise}[$\star\star$]\label{exo:pl:messaging-and-event-driven-architecture:5}
The \omterm{def:pl:position-and-pnl-services:service}{position service} keeps its state in memory and writes a snapshot every minute. Where should it keep its offset?
\end{exercise}

\begin{exercise}[$\star\star$]\label{exo:pl:messaging-and-event-driven-architecture:6}
Why does the outbox relay still produce repeats, and who removes them?
\end{exercise}

\begin{exercise}[$\star\star\star$]\label{exo:pl:messaging-and-event-driven-architecture:7}
\emph{Coding.} Make the at-least-once \omterm{def:pl:position-and-pnl-services:service}{position service} idempotent by storing applied fill identifiers, and run \texttt{schedules} on it.
\end{exercise}

\begin{exercise}[$\star\star\star$]\label{exo:pl:messaging-and-event-driven-architecture:8}
\emph{Find the flaw.} \textquotedbl{}We turned on exactly-once in the messaging system, so the positions cannot be double-counted any
more.\textquotedbl{}
\end{exercise}

\section{Problem: Double for Eleven Minutes}

\begin{problem}[{Weekend problem --- double for eleven minutes}]\label{pb:pl:messaging-and-event-driven-architecture:1}
The chapter's log, fills, crash schedules and trade-capture service.

\medskip\noindent\textbf{Part I --- The log.}
\begin{enumerate}
\item What does \omterm{def:pl:messaging-and-event-driven-architecture:pubsub}{publish--subscribe} decouple, and what does it cost?
\item When would you choose a broker, and when a brokerless transport?
\item How is a message addressed in a \omterm{def:pl:messaging-and-event-driven-architecture:log}{partitioned log}, and how is its partition chosen?
\item What does a \omterm{def:pl:messaging-and-event-driven-architecture:group}{consumer group} share, and what is a rebalance?
\item What does retention drop, and what must a new subscriber do if it has dropped events it needs?
\end{enumerate}

\medskip\noindent\textbf{Part II --- Crashes.}
\begin{enumerate}[resume]
\item What are the two steps of consuming a batch, and where can a crash fall?
\item How many crashes happened in each mode, and why fewer in one?
\item How many fills were lost or duplicated per 1\,000 crashes in each mode?
\item What is the worst end-of-day position error in each mode, and the median?
\item How does the atomic mode avoid both losses and duplicates?
\end{enumerate}

\medskip\noindent\textbf{Part III --- Producing.}
\begin{enumerate}[resume]
\item Why can the trade and its event not be written atomically without an outbox?
\item What does each naive design produce under crashes?
\item What does the outbox leave, and how is it removed?
\item What is \omterm{def:pl:messaging-and-event-driven-architecture:sourcing}{event sourcing}, and what does it require of the log?
\item What does a \omterm{def:pl:messaging-and-event-driven-architecture:dlq}{dead-letter queue} protect?
\end{enumerate}

\medskip\noindent\textbf{Part IV --- The verdict.}
\begin{enumerate}[resume]
\item State the \emph{named result}: duplicated and lost fills per thousand crashes under at-most-once, at-least-once and
idempotent-with-atomic-offset processing, and the maximum position error each produces.
\item Which mode should the \omterm{def:pl:position-and-pnl-services:service}{position service} use, and which a monitoring dashboard?
\item What does exactly-once in a messaging product cover?
\item How would you detect the hook's error within minutes?
\item In one sentence: where does exactly-once come from?
\end{enumerate}
\end{problem}

\section{Interview questions}

\begin{interviewq}[$\star$ \iqroles{developer}]\label{iq:pl:messaging-and-event-driven-architecture:1}
Explain at-most-once, at-least-once and exactly-once delivery.
\end{interviewq}

\begin{interviewq}[$\star\star$ \iqroles{developer}]\label{iq:pl:messaging-and-event-driven-architecture:2}
Your consumer applies fills to a database. How do you make sure a crash never double-counts a fill?
\end{interviewq}

\begin{interviewq}[$\star\star$ \iqroles{developer}]\label{iq:pl:messaging-and-event-driven-architecture:3}
How do you publish an event whenever a row is inserted in your database, without ever losing or inventing one?
\end{interviewq}

\begin{interviewq}[$\star\star$ \iqroles{developer}]\label{iq:pl:messaging-and-event-driven-architecture:4}
How do you keep the messages of one account in order when a topic has many partitions and many consumers?
\end{interviewq}

\begin{interviewq}[$\star\star\star$ \iqroles{developer}]\label{iq:pl:messaging-and-event-driven-architecture:5}
Design the event backbone that carries trades, fills and positions between the services of a trading firm.
\end{interviewq}
